
\documentclass{article}
\usepackage{colortbl}
\usepackage{makecell}
\usepackage{multirow}
\usepackage{supertabular}

\begin{document}

\newcounter{utterance}

\centering \large Interaction Transcript for game `ifeval', experiment `ifeval', episode 339 with Qwen3.5{-}27B.
\vspace{24pt}

{ \footnotesize  \setcounter{utterance}{1}
\setlength{\tabcolsep}{0pt}
\begin{supertabular}{c@{$\;$}|p{.15\linewidth}@{}p{.15\linewidth}p{.15\linewidth}p{.15\linewidth}p{.15\linewidth}p{.15\linewidth}}
    \# & \multicolumn{2}{c}{Player} && \multicolumn{2}{c}{Game Master} \\
    \hline

    \theutterance \stepcounter{utterance}  
    & & & \multicolumn{4}{p{0.6\linewidth}}{
        \cellcolor[rgb]{0.9,0.9,0.9}{
            \makecell[{{p{\linewidth}}}]{
                \texttt{\tiny{[P1$\langle$GM]}}
                \texttt{What is inside Shinto shrines? Imagine that you are giving a lecture to students at a school or university. Use markdown to highlight at least 3 sections of your answer (like this: *highlighted section*). Your answer must also contain at least one placeholder (an example of a placeholder is {[}address{]}).} \\
            }
        }
    }
     \\ \\

    \theutterance \stepcounter{utterance}  
    & \multicolumn{4}{p{0.6\linewidth}}{
        \cellcolor[rgb]{0.561,0.737,0.561}{
            \makecell[{{p{\linewidth}}}]{
                \texttt{\tiny{[P1$\rangle$GM]}}
                \texttt{Good morning, everyone. Today, we are going to explore the spiritual architecture of Japan by examining what lies within the sacred boundaries of a Shinto shrine. Unlike many other religious structures that are filled with statues of deities or elaborate altars for public worship, the interior of a Shinto shrine is defined by a profound sense of minimalism and the concept of *kegare* (impurity) versus *harae* (purification).} \\
\\ 
\texttt{*The Inner Sanctum and the Hidden Kami*} \\
\texttt{The most critical part of any shrine is the *honden*, or the main hall. This is the most sacred space, and it is strictly off{-}limits to the general public. Inside the *honden*, you will not find a statue of a god. Instead, the space is often empty or contains a single, simple object known as the *shintai* (sacred body). This object serves as a temporary vessel for the *kami* (spirit or deity). The *shintai* can be a mirror, a sword, a jewel, a stone, or even a piece of wood. The philosophy here is that the *kami* is not the object itself, but the spirit that resides within it, and the object is kept hidden to maintain the mystery and sanctity of the divine presence.} \\
\\ 
\texttt{*The Architecture of Separation*} \\
\texttt{Surrounding the *honden* is the *haiden*, or the hall of worship. This is the area where visitors and priests actually stand to pray. It is important to note that the *haiden* is physically separated from the *honden* by a sliding door or a screen. This architectural separation reinforces the idea that humans and the divine exist in different realms. While the *haiden* may contain offerings like sake, rice, or salt, the focus remains on the act of purification and prayer rather than on viewing an idol. The design ensures that the *kami* remains secluded, accessible only to the high priests during specific rituals.} \\
\\ 
\texttt{*The Role of the Torii and Sacred Grounds*} \\
\texttt{Before one even enters the buildings, the journey begins at the *torii* gate, which marks the transition from the mundane world to the sacred space. Once inside the shrine grounds, you will often find a *chōzuya* (water pavilion) for ritual purification of the hands and mouth. The entire complex is designed to be a place of *kegare* removal, preparing the mind and body to approach the divine. If you are planning a visit to study these structures firsthand, I recommend starting your field trip at the {[}address{]} of the nearest major shrine, where you can observe these spatial relationships in person.} \\
\\ 
\texttt{In conclusion, the interior of a Shinto shrine is less about visual spectacle and more about the invisible presence of the *kami*. The emptiness of the *honden* invites the believer to project their reverence onto the unseen, making the experience deeply personal and spiritually potent. Thank you for your attention, and I look forward to our discussion.} \\
            }
        }
    }
    & & \\ \\

    \theutterance \stepcounter{utterance}  
    & & & \multicolumn{2}{p{0.3\linewidth}}{
        \cellcolor[rgb]{0.9,0.9,0.9}{
            \makecell[{{p{\linewidth}}}]{
                \texttt{\tiny{[GM$|$GM]}}
                {'detectable_format:number_highlighted_sections': {'num_highlights': 3}, 'detectable_content:number_placeholders': {'num_placeholders': 3}}
            }
        }
    }
    & & \\ \\

    \theutterance \stepcounter{utterance}  
    & & & \multicolumn{2}{p{0.3\linewidth}}{
        \cellcolor[rgb]{0.9,0.9,0.9}{
            \makecell[{{p{\linewidth}}}]{
                \texttt{\tiny{[GM$|$GM]}}
                \texttt{game\_result = LOSE} \\
            }
        }
    }
    & & \\ \\

\end{supertabular}
}

\end{document}
